\chapter{Background}
\label{chap:background}

\section{Markov Decision Processes}

\begin{definition}[Finite-Horizon MDP]
	A finite-horizon Markov Decision Process is a tuple $(\mathcal{S}, \mathcal{A}, P, R, T)$ where:
	\begin{itemize}
		\item $\mathcal{S}$ is a finite state space
		\item $\mathcal{A}$ is a finite action space
		\item $P: \mathcal{S} \times \mathcal{A} \times \mathcal{S} \to [0,1]$ is the transition probability function
		\item $R: \mathcal{S} \times \mathcal{A} \to \R$ is the reward function
		\item $T \in \mathbb{N}$ is the time horizon
	\end{itemize}
\end{definition}

A policy $\pi: \mathcal{S} \times \{0, \ldots, T-1\} \to \mathcal{A}$ maps states and time to actions. The value function under policy $\pi$ is:
%
\begin{equation}
	V^\pi_t(s) = \E^\pi\left[\sum_{k=t}^{T-1} R(s_k, a_k) \mid s_t = s\right]
\end{equation}

The optimal value function satisfies the Bellman optimality equation:
%
\begin{equation}
	V^*_t(s) = \max_{a \in \mathcal{A}} \left\{ R(s,a) + \sum_{s' \in \mathcal{S}} P(s'|s,a) V^*_{t+1}(s') \right\}
\end{equation}

\section{Approximate Dynamic Programming}

When the state space is large, exact dynamic programming becomes computationally intractable. Approximate Dynamic Programming addresses this through value function approximation.

\subsection{Linear Value Function Approximation}

We approximate the value function as a linear combination of basis functions:
%
\begin{equation}
	\hat{V}(s; \theta) = \sum_{i=1}^k \theta_i \phi_i(s) = \phi(s)^\top \theta
\end{equation}
%
where $\phi: \mathcal{S} \to \R^k$ maps states to feature vectors and $\theta \in \R^k$ are learnable weights.

\subsection{Temporal Difference Learning}

The TD(0) update for linear VFA is:
%
\begin{equation}
	\theta \leftarrow \theta + \alpha \left[r + \hat{V}(s'; \theta) - \hat{V}(s; \theta)\right] \phi(s)
\end{equation}
%
where $\alpha > 0$ is the learning rate.

\section{Proto-Value Functions}

Proto-Value Functions \citep{mahadevan2007proto} provide a task-independent basis for value function approximation derived from the topology of the state space.

\subsection{Graph Laplacian}

Given a weighted graph $G = (V, E, W)$ representing state transitions:
%
\begin{equation}
	L = D - W
\end{equation}
%
where $W$ is the weighted adjacency matrix and $D$ is the degree matrix with $D_{ii} = \sum_j W_{ij}$.

The normalized Laplacian is:
%
\begin{equation}
	\mathcal{L} = I - D^{-1/2} W D^{-1/2}
\end{equation}

\subsection{Spectral Basis Functions}

The eigenvectors of $\mathcal{L}$ corresponding to the smallest eigenvalues form a smooth basis over the state space. Let $0 = \lambda_1 \leq \lambda_2 \leq \cdots \leq \lambda_{|\mathcal{S}|}$ be the eigenvalues with eigenvectors $v_1, \ldots, v_{|\mathcal{S}|}$.

The first $k$ non-trivial eigenvectors $\{v_2, \ldots, v_{k+1}\}$ serve as PVFs:
%
\begin{equation}
	\phi_i(s) = v_{i+1}(s), \quad i = 1, \ldots, k
\end{equation}

\subsection{Interpretation}

PVFs capture the "geometry" of the state space:
\begin{itemize}
	\item Lower eigenvalues correspond to smoother functions over the graph
	\item They identify bottlenecks and clusters in state-transition structure
	\item Task-independent: derived from dynamics, not rewards
\end{itemize}

\section{Successor Features}

Successor features \citep{dayan1993improving, barreto2017successor} represent states by their expected future occupancy.

\begin{definition}[Successor Representation]
	The successor representation under policy $\pi$ with discount factor $\gamma$ is:
	%
	\begin{equation}
		\psi^\pi(s) = \E^\pi\left[\sum_{t=0}^\infty \gamma^t \mathbb{I}[s_t = s] \mid s_0 = s\right]
	\end{equation}
\end{definition}

Successor features generalize this to feature expectations:
%
\begin{equation}
	\psi^\pi(s) = \E^\pi\left[\sum_{t=0}^\infty \gamma^t \phi(s_t) \mid s_0 = s\right]
\end{equation}

These can be learned via TD updates and enable transfer learning across tasks with different reward functions.

\section{Related Work}

[Summary of relevant literature on airport scheduling, ADP applications in operations research, spectral methods in RL, etc.]